\SourceLeaf{85}
\fancyhead[L]{\small Intensitas vis magneticae terrestris}
\renewcommand{\Article}[1]{\par\medskip\Needspace{4\baselineskip}\pdfbookmark[1]{Art. #1}{magnetism-art-#1}\centerline{\textbf{#1.}}\nobreak\smallskip}
\pdfbookmark[0]{Intensitas vis magneticae terrestris ad mensuram absolutam revocata}{magnetism-body}
\par\medskip
\centerline{INTENSITAS}
\centerline{VIS MAGNETICAE TERRESTRIS}
\centerline{AD MENSURAM ABSOLUTAM REVOCATA.}
\medskip
Ad determinationem completam vis magneticae telluris in loco dato tria elementa requiruntur: declinatio seu angulus inter planum in quo agit atque planum meridianum; inclinatio directionis ad planum horizontale; denique tertio loco intensitas. Declinatio, quae respectu omnium applicationum ad usus nauticos atque geodaeticos tamquam elementum primarium consideranda est, statim ab initio observatores atque physicos exercuit, qui etiam inclinationi curas assiduas iam per saeculum dicaverunt. Contra elementum tertium, intensitas, licet aeque dignum scientiae obiectum, usque ad tempora recentiora penitus neglectum mansit. Illustri \textsc{Humboldt} inter tot alias ea quoque laus debetur, quod primus fere huic argumento animum advertit, inque itineribus suis magnam copiam observationum circa intensitatem relativam magnetismi terrestris congessit, e quibus continuum incrementum huius intensitatis, dum ab aequatore magnetico versus polum progredimur, innotuit. Permulti physici, vestigiis huius naturae scrutatoris insistentes, iam tantam determinationum copiam contulerunt, ut clarissimus et de magnetismo terrestri meritissimus vir, \textsc{Hansteen}, specimen mappae universalis lineas isodynamicas exhibentis nuper iam edere potuerit.

Methodus, qua in hoc negotio utuntur, consistit in observatione temporis, per quod eadem acus magnetica in locis diversis numerum oscillationum eundem perficit, seu numeri oscillationum in temporis intervallo eodem, intensitasque quadrato numeri oscillationum in tempore dato proportionalis ponitur: hoc modo inter se comparantur intensitates totales, dum acus inclinatoria in centro gravitatis suspensa circa axem horizontalem ad meridianum magneticum normalem oscillat, seu intensitates vis horizontalis, dum acus horizontalis circa axem ver\SourceLeaf{86}ticalem vibratur: posterior observandi modus maiorem praecisionem fert, et quae inde resultant, cognitis inclinationibus, facile ad intensitates totales reducuntur.

Manifesto nervus huius methodi pendet a suppositione, distributionem magnetismi liberi in particulis acus ad talem comparationem adhibitae in singulis experimentis invariatam mansisse: si enim vis magnetica acus lapsu temporis aliquantam debilitationem passa esset, ob id ipsum postea tardius oscillaret, observatorque talis mutationis inscius intensitati magnetismi terrestris pro loco posteriori valorem nimis parvum tribueret. Quodsi experimenta temporis intervallum mediocre complectuntur, acusque ex chalybe bene durato confecta et magnetismo caute imbuta in usum vocatur, considerabilis vigoris debilitatio parum utique metuenda erit; praeterea incertitudo minuetur, si plures acus ad comparationes adhibentur; denique isti suppositioni maior fides accrescet, si peracto itinere acus in loco primo tempus vibrationis non mutasse invenitur. Sed quaecunque cautelae adhibeantur, lenta aliqua debilitatio vis acus vix evitari, adeoque talis consensus post longiorem absentiam raro exspectari poterit; quapropter in comparatione intensitatum pro locis terrae valde dissitis plerumque tantam praecisionem et certitudinem, quantam desiderare debemus, attingere haud licebit.

Ceterum hoc methodi incommodum minus grave est, quamdiu tantum de comparatione intensitatum simultanearum vel temporibus non longe inter se distantibus respondentium agitur. At quum experientia docuerit, tum declinationem tum inclinationem in loco dato mutationes continuas pati, quae post multos annos pergrandes evadant, dubium esse nequit, quin intensitas quoque magnetismi terrestris similibus mutationibus quasi saecularibus obnoxia sit; manifesto autem, quatenus de \textit{hac} quaestione agitur, methodus ista prorsus inefficax evadit. Et tamen, ad scientiae naturalis incrementum summopere desiderandum foret, ut haec ipsa quaestio gravissima in plenissimam lucem promoveatur, quod certo fieri nequit, nisi methodo pure comparativa abrogata alia substituitur, quae a fortuitis acuum inaequalitatibus prorsus independens intensitatem magnetismi terrestris ad unitates stabiles mensurasque absolutas revocat.

Haud difficile est, principia theoretica stabilire, quibus talis methodus, diu iam in votis habita, inniti debet. Multitudo oscillationum, quas acus in tempore dato perficit, pendet tum ab intensitate magnetismi terrestris, tum a constitutione acus, puta a momento statico elementorum magnetismi liberi in illa contentorum, atque ab eius momento inertiae. Quum hoc momentum inertiae absque difficul\SourceLeaf{87}tate assignari possit, patet, observationem oscillationum nobis suppeditare productum ex intensitate magnetismi terrestris in momentum staticum magnetismi acus: sed hae duae quantitates separari nequeunt, nisi observationibus alius generis accitis, quae diversam earum combinationem implicant. Ad hunc finem accedat acus secunda, quae exponatur actioni et magnetismi terrestris et magnetismi acus primae, ut, quam rationem inter se teneant hae duae actiones, explorari possit. Utraque quidem actio pendebit a distributione magnetismi liberi in secunda acu: sed posterior insuper a constitutione acus primae, distantia centrorum, positione rectae centra iungentis respectu axium magneticorum utriusque acus, denique a lege, quam attractiones et repulsiones magneticae sequuntur. Immortalis \textsc{Tobias Mayer} primus iam coniectaverat, hanc legem cum lege gravitationis eatenus convenire, quod illae quoque actiones decrescant in ratione duplicata distantiarum: experimenta clarissimorum virorum \textsc{Coulomb} et \textsc{Hansteen} magnam huic coniecturae plausibilitatem conciliaverunt, experimentaque novissima eam ultra omne dubium evehunt. Sed probe attendendum est, hanc legem referri ad singula elementa magnetismi liberi: effectus totalis corporis magnetismo imbuti longe aliter se habebit, atque in distantiis permagnis, uti ex illa ipsa lege deducere licet, proxime ad rationem inversam triplicatam distantiarum accedet, ita ut actio acus per cubum distantiae multiplicata, distantia ceteris paribus continuo crescente, ad valorem constantem asymptotice convergat, qui, dum distantiae, linea arbitraria pro unitate accepta, per numeros exprimuntur, cum actione vis terrestris homogeneus atque comparabilis erit. Per idoneam experimentorum adornationem et tractationem limes huius rationis eruendus est; qui quum tantummodo momentum staticum magnetismi primae acus involvat, iam habebitur quotiens e divisione huius momenti per intensitatem vis terrestris ortus, qui comparatus cum producto harum quantitatum antea eruto eliminationi istius momenti statici inserviet, atque valorem intensitatis magnetismi terrestris suppeditabit.

Quod attinet ad modum, actiones magnetismi terrestris et acus primae in acum secundam ad experimenta revocandi, duplex via patet, quum acum secundam vel in statu motus vel in statu aequilibrii observare possimus. Prior modus eo redit, ut oscillationes huius acus observentur, dum actio magnetismi terrestris coniungitur cum actione acus primae in distantia idonea ita collocatae, ut ipsius axis sit in meridiano magnetico per centrum acus oscillantis ducto: hoc pacto oscillationes vel accelerabuntur vel retardabuntur, prout poli amici vel inimici sibi \SourceLeaf{88}mutuo obversi sunt, comparatioque vel temporum vibrationum pro utraque acus primae positione inter se, vel temporis alterutrius cum tempore vibrationum sub sola magnetismi terrestris actione (remota acu prima) locum habentium, rationem huius vis ad actionem primae acus docebit. Alter modus acum primam ita collocat, ut directio vis quam in regione acus secundae libere suspensae exercet, faciat angulum (e. g. rectum) cum meridiano magnetico, quo pacto haec ipsa a meridiano magnetico deflectetur, et e magnitudine deflexus ratio inter vim magneticam terrestrem atque actionem acus primae concludetur.

Ceterum modus prior essentialiter convenit cum eo, quem ill. \textsc{Poisson} iam ante aliquot annos proposuit. Sed experimenta ad ipsius normam a nonnullis physicis tentata, quae quidem mihi innotuerunt, vel successu omnino caruerunt, vel rudem tantummodo approximationem praebuerunt.

Difficultas rei inde potissimum pendet. quod ex actionibus acus in distantiis mediocribus observatis computari debet limes aliquis, qui ad distantiam quasi infinite magnam refertur, et quod eliminationes ad hunc finem necessariae tanto magis a levissimis observationum erroribus turbantur, imo pervertuntur, quo plures incognitae a statu acuum individuali pendentes eliminandae sunt: ad multitudinem perparvam incognitarum autem res tunc tantummodo deduci potest, ubi actiones in distantiis satis magnis (respectu longitudinis acuum) observatae sunt, adeoque ipsae iam perparvae evaserunt. Sed ad actiones tam parvas accurate mensurandas subsidia practica hactenus usitata non sufficiunt.

Ante omnia itaque in id mihi incumbendum esse vidi, ut subsidia nova parem, per quae tum tempora oscillationum, tum directiones acuum longe maiori praecisione, quam hactenus licuit, observare ac metiri possem. Labores ad hunc finem suscepti et per plures menses continuati, in quibus a praestantissimo \textsc{Weber} multifariam adiutus sum, ad scopum exoptatum ita perduxerunt, ut exspectationem non modo non fefellerint, sed longe superaverint, nec iam quidquam desiderandum restet, ad praecisionem experimentorum subtilitati observationum astronomicarum aequiparandam, nisi locus ab influxu ferri propinqui atque agitationum aëris plene securus. Adsunt duo apparatus simplicitate non minus quam praecisione quam praebent insignes, quorum descriptionem quidem ad aliam occasionem mihi reservare debeo, dum experimenta ad determinandam intensitatem magnetismi terrestris hactenus in observatorio nostro instituta physicis in hac commentatione trado.
